%------------------------------------------------------------------------------%
% LINEAR OPERATOR EQUATIONS                                                    %
%------------------------------------------------------------------------------%
% File  : Linear_Operator_Equations.tex                                        %
% Author: Klaus Hoeflinger, Beethovenstrasse 11, D-73117 Wangen                %
% Email : khoeflinger@t-online.de                                              %
%------------------------------------------------------------------------------%

%------------------------------------------------------------------------------%
%                       DOCUMENT CLASS and INCLUDE FILES                       %
%------------------------------------------------------------------------------%
\documentclass[11pt,twoside]{article}
\usepackage{geometry}
\geometry{paperheight=241mm,
          paperwidth=152mm,
          top=22mm,
          bottom=17mm,
          left=20mm,
          right=20mm,
          textwidth=112mm,
          textheight=202mm,
          headheight=10mm,
          headsep=3mm,
          footskip=9mm,
          includehead,
          includefoot,
          marginparsep=0mm,
          marginparwidth=0mm}

\renewcommand{\baselinestretch}{0.945}\normalsize

\AtBeginDocument{\setlength\abovedisplayskip{2mm}}
\AtBeginDocument{\setlength\belowdisplayskip{2mm}}

%------------------------------------------------------------------------------%
% define title formats                                                         %
%------------------------------------------------------------------------------%
\usepackage{titlesec}

\renewcommand{\thesection}{\arabic{section}}
\renewcommand{\sfdefault}{FiraSans}
\renewcommand{\ttdefault}{pcr}
\renewcommand{\rmdefault}{antiqua}

\titleformat{\part}[display]
{\normalfont \large \rmfamily \bfseries \centering}
{\small{\thepart.}}{2pt}{}

\titleformat{\section}[runin]
{\normalfont \rmfamily \bfseries}
{\thesection}{1pt}{.\;}[.]

%\titleformat{\section}
%{\normalfont \bfseries \small \filcenter}
%{\thesection.\;}{0mm}{}

\titleformat{\subsection}[runin]
{\normalfont \itshape}
{}{0mm}{}

\titleformat{\subsubsection}[runin]
{\normalfont \itshape}
{}{}{}{}

\titleformat{\paragraph}[runin]
{\normalfont}
{}{}{}{}

\titleformat{\subparagraph}[runin]
{\normalfont}
{}{}{}{}

\titlespacing*{\part}          { 0pt}{ 0pt}{ 8pt}
\titlespacing*{\section}       { 0pt}{ 7pt}{ 4pt}
\titlespacing*{\subsection}    { 0pt}{14pt}{ 6pt}
\titlespacing*{\subsubsection} { 0pt}{ 6pt}{12pt}
\titlespacing*{\paragraph}     { 0pt}{ 6pt}{12pt}
\titlespacing*{\subparagraph}  { 0pt}{ 0pt}{12pt}

%------------------------------------------------------------------------------%
% define enumerate parameters                                                  %
%------------------------------------------------------------------------------%
\usepackage{enumitem}
\setlength{\parindent}{3pt}
\setlength{\parskip}  {0pt}

%\setlist{noitemsep}

\setlist[enumerate]{
         leftmargin=12pt,
         rightmargin=0pt,
         itemsep=1pt,
         itemindent=3pt,
         topsep=3pt,
         labelsep=4pt,
         partopsep=0pt,
         labelwidth=0pt,
         listparindent=0pt,
         parsep=0pt}

\setlist[itemize]{
         leftmargin=12pt,
         rightmargin=0pt,
         itemsep=0pt,
         itemindent=1pt,
         topsep=1pt,
         labelsep=4pt,
         partopsep=1pt,
         labelwidth=0pt,
         listparindent=0pt,
         parsep=0pt}

%------------------------------------------------------------------------------%
% define packages                                                              %
%------------------------------------------------------------------------------%
\usepackage{upref}
\usepackage{amssymb}
\usepackage{slantsc}
\usepackage{amsmath}
\usepackage{amsthm}
\usepackage{thmtools}
\usepackage{amscd}
\usepackage{verbatim}
\usepackage{latexsym}
\usepackage{multicol}
\usepackage{graphicx}
\usepackage{setspace}
\usepackage[ngerman,english]{babel}
\usepackage[protrusion=true,expansion=true]{microtype}
\usepackage{tikz}
\usetikzlibrary{arrows,arrows.meta,decorations.markings}
\usepackage{mathrsfs}
\usepackage{amsfonts}
\usepackage{datetime2}
\usepackage{textgreek}
\usepackage{hyperref}
\usepackage{pifont}
\usepackage{relsize}
\usepackage[UKenglish]{isodate}
\usepackage[skip=0mm]{parskip}
\usetikzlibrary{decorations.pathmorphing}
\relsize{-0.05}

\usepackage{mlmodern}
\usepackage[T5,T1]{fontenc}
\usepackage{ragged2e}

%\usepackage{mathabx}
%\usepackage{times}

%\usepackage{fontspec}
%\setmainfont{Nimbus Roman No9 L}
%\usepackage[T1]{fontenc}

%------------------------------------------------------------------------------%
% define page layout                                                           %
%------------------------------------------------------------------------------%
\usepackage{fancyhdr}
\pagestyle{fancy}
\setlength{\headheight}{30.0pt}
\renewcommand{\headrulewidth}{0pt}

\AtBeginDocument{
  \addtolength\abovedisplayskip{-0.0\baselineskip}
  \addtolength\belowdisplayskip{-0.0\baselineskip}
}

\fancyhead{}
\fancyfoot{}
\addtolength{\headwidth}{0.02\marginparwidth}

\makeatletter
\let\ps@plain\ps@fancy
\makeatother

\renewcommand\sectionmark[1]
{
 \def\sectionname{#1}
 \markright{\thesection #1}
}

\makeatletter
\@addtoreset{section}{part}
\makeatother

%------------------------------------------------------------------------------%
% define footnote without number                                               %
%------------------------------------------------------------------------------%
\newcommand{\myfootnote}[1]{%
\renewcommand{\thefootnote}{}%
\footnotetext{#1}%
\renewcommand{\thefootnote}{\arabic{footnote}}%
}

%------------------------------------------------------------------------------%
% define numbering in equations                                                %
%------------------------------------------------------------------------------%
\numberwithin{equation}{section}

%------------------------------------------------------------------------------%
% define newtheorem                                                            %
%------------------------------------------------------------------------------%
\declaretheoremstyle[
headfont=\scshape, 
headindent = 0mm, %\parindent,
%postheadhook = {\hspace{0mm}\newline},
%postheadspace = \newline, 
spaceabove = 3mm, 
spacebelow = 3mm,
numberwithin=section,
name=Theorem]{khMATH}
\declaretheorem[style=khMATH]{theorem}

\declaretheoremstyle[
headfont=\bfseries, 
headindent = 0mm, %\parindent,
%postheadhook = {\hspace{0mm}\newline},
%postheadspace = \newline, 
spaceabove = 3mm, 
spacebelow = 3mm,
numberwithin=part,
name=Theorem]{khMATH1}
\declaretheorem[style=khMATH1]{theorem1}

\declaretheoremstyle[
headfont=\scshape, 
headindent = 0mm, %\parindent,
%postheadhook = {\hspace{0mm}\newline},
%postheadspace = \newline, 
spaceabove = 3mm, 
spacebelow = 3mm,
numberwithin=section,
name=Corollary]{khMATH1}
\declaretheorem[style=khMATH1]{corollary}

%            name         counter     label              numeration-mode
\theoremstyle{plain}
\newtheorem*{introduction}            {\sc Introduction}
\newtheorem {standard}                {}                 [section]
\newtheorem*{definition}              {\sc Definition}
\newtheorem{satz1}                    {\sc Satz}
\newtheorem{lemma}                    {\sc Lemma}
\newtheorem*{proposition}             {\sc Proposition}
%\newtheorem{corollary}               {\sc Corollary}
\newtheorem*{corollar}                {\sc Korollar}
\newtheorem*{remark}                  {\sc Remark}
\newtheorem*{remarks}                 {\sc Remarks}
\newtheorem*{example}                 {\sc Example}
\newtheorem*{examples}                {\sc Examples}
\newtheorem*{theo}                    {\sc Theorem}

%------------------------------------------------------------------------------%
% define tabular                                                               %
%------------------------------------------------------------------------------%
\usepackage{tabularx}
\newcolumntype{L}[1]{>{\raggedright\arraybackslash}p{#1}}
\newcolumntype{C}[1]{>{\centering\arraybackslash}  p{#1}} 
\newcolumntype{R}[1]{>{\raggedleft\arraybackslash} p{#1}} 

%------------------------------------------------------------------------------%
% define \sh, \zB                                                              %
%------------------------------------------------------------------------------%
\def\dsh{{d.\,h.}}
\def\ie{{i.\,e.}}
\def\zB{z.\,B.}
\renewcommand{\abstractname}{ }

%------------------------------------------------------------------------------%
% change default spacing for cases                                             %
%------------------------------------------------------------------------------%
\makeatletter
\renewcommand*\env@cases[1][1.6]{%
  \let\@ifnextchar\new@ifnextchar
  \left\lbrace
  \def\arraystretch{#1}%
  \array{@{}l@{\quad}l@{}}%
}
\makeatother

\newenvironment{rcases}
  {\left.\begin{aligned}}
  {\end{aligned}\right\rbrace}

%------------------------------------------------------------------------------%
% change default for \predisplaypenalty                                        %
%------------------------------------------------------------------------------%
\predisplaypenalty=990
\allowdisplaybreaks \numberwithin{equation}{section}

%------------------------------------------------------------------------------%
% general notations                                                            %
%------------------------------------------------------------------------------%
\def\*{\star}
\def\={\,=\,}
\def\+{\,+\,}
\def\xsubset{{\,\subset\,}}
\def\xtimes{{\,\times\,}}
\def\xeof{{\,\in\,}}
\def\eq{{\,=\,}}
\def\xto{{\,\to\,}}
\def\eqdef{\,:=\,}
\def\:{{\,:\,}}
\def\ele{{\,\in\,}}
\def\exists{\mbox{ exists }}
\def\and{\mbox{ and }}
\def\for{\mbox{ for }}
\def\fa{\mbox{ for all }}
\def\fe{\mbox{ for each }}
\def\qfa{\quad \mbox{ for all }}
\def\df{\,\dotfill}
\def\iff{if and only if }
\def\wrt{with respect to}

\makeatletter
\newcommand \Df {\leavevmode \cleaders \hb@xt@ .70em{\hss .\hss }\hfill \kern \z@}
\makeatother

\def\sql{\mathfrak{a}}               % sesquilinear form a
\def\DF{B}                           % Dirichlet form a
%------------------------------------------------------------------------------%
% number sets and special elements                                             %
%------------------------------------------------------------------------------%
\def\P{\mathbb{H}}                   % half plane of $\C$
\def\J{I}                            % interval I
\def\N{{\mathbb{N}}}                 % unsigned integers
\def\Q{{\mathbb{Q}}}                 % rational integers
\def\Z{{\mathbb{Z}}}                 % integers
\def\R{{\mathbb{R}}}                 % real numbers
\def\Rn{{\mathbb{R}}^n}              % real numbers in Rn
\def\Rd{{\mathbb{R}}^d}              % real numbers in Rn
\def\C{{\mathbb{C}}}                 % complex numbers
\def\F{{\mathbb{K}}}                 % arbitrary field
\def\H{{\mathbb{H}}}                 % half of the complex plane
\def\Torus{{\mathbb{T}}}             % complex circle line
\def\T{\tau}                         % upper bound of interval (0,t)
\def\sign#1{\operatorname{sign}(#1)} % sign of a number

%------------------------------------------------------------------------------%
% set theory                                                                   %
%------------------------------------------------------------------------------%
\def\G{G}                            % domain $G$
\def\clG{\overline{\G}}              % closure of $G$
\def\bndG{\partial \G}               % boundary of $G$
\def\setA{\mathcal{A}}               % set A
\def\setB{\mathcal{B}}               % set B
\def\setC{\mathcal{C}}               % set C
\def\setM{M}                         % set M
\def\setN{N}                         % set N
\def\setS{\mathcal{S}}               % set S
\def\famF{\mathfrak{F}}              % family of sets
\def\indI{\mathcal{I}}               % index set I
\def\pot#1{\mathfrak{P}(#1)}         % power set

\def\relR{\mathbf{R}}                % relation R
\def\mapf{f}                         % mapping f
\def\mapg{g}                         % mapping g

\def\set#1#2{\{ \, #1 \,: \, #2 \, \}}
\def\tensor#1#2{#1 \otimes #2}
\def\Tens{\oplus}

\def\cal#1{{\mathcal{#1}}}
\def\aut#1{\operatorname{Aut}(#1)}

\def\cross#1#2{#1 {\,\times\,} #2}
\def\multicross#1#2{#1 \times  \ldots \times  #2}
\def\multitens#1#2 {#1 \otimes \ldots \otimes #2}

%------------------------------------------------------------------------------%
% group theory                                                                 %
%------------------------------------------------------------------------------%
\def\1{{\mathsf{1}}}                % unit element of a monoid
\def\0{{\mathsf{0}}}                % unit element of an Abelian monoid
\def\kernel#1{\mathfrak{N}(#1)}     % kernel of a homomorphism
\def\cokernel#1{\mathfrak{C}{(#1)}} % cokernel of a homomorphism
\def\Hom#1{\operatorname{Hom}(#1)}

%------------------------------------------------------------------------------%
% linear algebra                                                               %
%------------------------------------------------------------------------------%
\def\vecV{\mathit{V}}               % vector space V
\def\vecH{\mathit{H}}               % vector space H
\def\vecW{\mathit{W}}               % vector space W
\def\vecU{\mathit{U}}               % subspace U
\def\convC{\mathcal{C}}             % convex set C
\def\basH{\mathfrak{B}}             % Hamel base B

\def\LO#1{\mathscr{L}(#1) }         % algebra of linear transformations

\def\scalar#1#2{ \langle #1,#2 \rangle }
\def\trace#1{\operatorname{tr}(#1)}
\def\dim#1{{\operatorname{dim} \, #1 }}
\def\codim#1{{\operatorname{codim} \, #1 }}
\def\dim{\operatorname{dim}}
\def\codim{\operatorname{codim}}
\def\ind{\operatorname{ind}}

\def\genG{\cal{G}}

\def\algA{\cal{A}}
\def\algB{\cal{B}}
\def\algU{\cal{U}}

\def\idI{\cal{I}}

%------------------------------------------------------------------------------%
% topology                                                                     %
%------------------------------------------------------------------------------%
\def\compK{{K}}          % compact set C
\def\openO{\mathcal{O}}  % open set O
\def\openU{\mathcal{U}}  % open set U
\def\U{\mathcal{U}}      % neighbourhood U
\def\V{\mathcal{V}}      % neighbourhood V
\def\d{\mathit{d}}       % metric function d

\def\interior#1{\mathring{#1}}
\def\closure#1{{\overline{#1}}}

\def\topT{\mathcal{T}}   % topology T
\def\topX{X}             % topological space X
\def\topY{Y}             % topological space Y
\def\topZ{Z}             % topological space Z
\def\metrS{M}            % metric space M

\def\grpG{G}             % topological group G
\def\grpA{A}             % topological group A
\def\grpB{B}             % topological group B
\def\grpC{C}             % topological group C
\def\grpH{H}             % topological group H
\def\grpU{U}             % topological group U
\def\grpV{V}             % topological group V
\def\topG{G}             % topological group G
\def\topH{H}             % topological group H
\def\topU{U}             % topological subgroup U

\def\subS{\cal{S}}
\def\riR{R}              % ring R

%------------------------------------------------------------------------------%
% calculus                                                                     %
%------------------------------------------------------------------------------%
\def\supp#1{\operatorname{supp}(#1)}
\def\singsupp#1{\operatorname{sing \, supp}(#1)}

\def\re{\operatorname{Re}}
\def\im{\operatorname{Im}}
\def\grad{\operatorname{grad}}

%\def\abs#1 {\left|#1\right|}
\def\abs#1 {\lvert #1 \rvert}
%\def\abs#1 {\left|\,#1\,\right|}

%------------------------------------------------------------------------------%
% measure theory                                                               %
%------------------------------------------------------------------------------%
\def\salgA{\Sigma}               % sigma-algebra S
\def\salgB{\cal{B}}              % sigma-algebra B
\def\Borel#1{\cal{B}_{#1}}       % Borel-algebra 
\def\LB{\lambda}                 % Lebesgue-Borel measure

%------------------------------------------------------------------------------%
% function spaces                                                              %
%------------------------------------------------------------------------------%
\def\Lp{L^p(\G)}                 % spaces L_p
\def\Lq{L^q(\G)}                 % spaces L_q
\def\Lh{L^2(\G)}                 % spaces L_2
\def\Li{L^{\infty}(\G)}          % spaces L_infty
\def\Lc{L_{loc}^1(\G)}           % spaces L_1^{loc}
\def\Ck{C^k(\G)}                 % spaces C_k

\def\BV#1{BV(#1) }               % set of functions of bounded variation
\def\AC#1{AC(#1) }               % set of absolutely continuous functions

\def\MP#1{\mathscr{M}^+(#1) }    % set of positive measures
\def\MS#1{\mathscr{M}^-(#1) }    % vector space of finite signed measures
\def\MC#1{\mathscr{M}(#1) }      % vector space of complex measures
\def\MCR#1{\mathscr{M}_{r}(#1) } % vector space of regular complex measures

\def\SobW{W^{k,p}(\G)}                % Sobolev space W
\def\SobO{\mathring{W}^{k,p}(\G)}     % Sobolev space W
%\def\WN#1#2{\mathring{W}^{#1}(#2) }   % Sobolev space W
\def\WN#1#2{W_0^{#1}(#2) }   % Sobolev space W
%\def\HN#1#2{\mathring{H}^{#1}(#2) }   % Sobolev space H
\def\HN#1#2{H_0^{#1}(#2) }   % Sobolev space H
%------------------------------------------------------------------------------%
% functional analysis                                                          %
%------------------------------------------------------------------------------%
\def\snorm#1{\lVert #1 \rVert}
\newcommand{\norm}[1]{\left\lVert #1 \right\rVert}

\def\topV{V}                       % topological vector space V
\def\topW{W}                       % topological vector space W
\def\normS{X}                      % normed space

\def\banX{\mathit{X}}              % Banach space X
\def\banY{\mathit{Y}}              % Banach space Y
\def\banZ{\mathit{Z}}              % Banach space Z
\def\dbanX{\mathit{X^*}}           % dual space of Banach space X
\def\dbanY{\mathit{Y^*}}           % dual space of Banach space Y
\def\dbanZ{\mathit{Z^*}}           % dual space of Banach space Z

\def\banA{\mathit{A}}              % Banach algebra A
\def\banB{\mathit{B}}              % Banach algebra B
\def\banC{\mathit{C}}              % C*-algebra C


\def\domain#1{\mathfrak{D}(#1)}    % domain
\def\range#1{\mathfrak{R}(#1)}     % range

\def\isoJ{\mathfrak{J}}            % isometry J
\def\repA{{\cal{A}_{\sql}}}        % representation operator A
\def\linS{{S}}                     % linear operator S
\def\linG{{G}}                     % linear operator G
\def\linL{{L}}                     % linear operator L
\def\linT{{T}}                     % linear operator A
\def\linA{\mathit{A}}              % linear operator A
\def\linB{{B}}                     % linear operator B
\def\linM{{M}}                     % linear operator M
\def\linU{{U}}                     % linear operator U
\def\linI{{I}}                     % linear operator I
\def\A{\cal{A}}                    % linear differential operator A
\def\BDO{\cal{B}}                  % linear boundary differential operator B
\def\linU{{U}}                     % unitary operator U
\def\linP{{P}}                     % projection P
\def\compA{k}                      % compact operator k
\def\linFR{f}                      % finite rank operator f
\def\fredA{F}                      % Fredholm operator F

\def\HAM{\cal{H}}                  % Hamilton operator
\def\PA{\partial^{\alpha}}         % elementary differential operator
\def\DO{\cal{A}}                   % linear differential operator
\def\BC{\cal{B}}                   % boundary operator
\def\SLO{\cal{L}_{p,q}}            % Sturm-Liouville operator
\def\MO{\cal{L}_{M}}               % momentum operator

\def\HS#1{S(#1)           }        % algebra of Hilbert-Schmidt operators
\def\FR#1{F(#1)           }        % algebra of finite rank operators
\def\TC#1{N(#1)           }        % algebra of trace class operators
\def\BU#1{\mathscr{U}(#1) }        % algebra of unitary linear operators
\def\BO#1{\mathscr{L}(#1) }        % algebra of bounded linear operators
\def\CO#1{\mathscr{C}(#1) }        % algebra of compact linear operators
\def\CL#1{\mathscr{C}(#1) }        % algebra of compact linear operators
\def\FO#1{\Phi(#1) }               % set of Fredholm operators
\def\FK#1#2{\Phi_{#2}(#1) }        % set of Fredholm operators with index k

\def\hilH{\mathit{H}}              % Hilbert space H
\def\clU{\mathit{U}}               % closed subspace U

\def\orthO{\mathfrak{O}}           % orthonormal system O
\def\orthS{\mathfrak{S}}           % orthonormal system S
\def\basB{\mathfrak{B}}            % basis B of a Hilbert space

\def\GL#1 {\mathbf{GL}(#1)}        % linear group
\def\OG#1 {\mathbf{O}(#1)}         % linear group
\def\SOG#1 {\mathbf{SO}(#1)}       % linear group

%------------------------------------------------------------------------------%
% semigroup theory                                                             %
%------------------------------------------------------------------------------%
\def\SG{$\{S(t)\}_{t \geq 0}$}
\def\TG{$\{T(t)\}_{t \geq 0}$}
\def\GG{$\{G(t)\}_{t \geq 0}$}
\def\EG{$\{\cal{E}(t)\}_{t \geq 0}$}

%------------------------------------------------------------------------------%
% distribution theory                                                          %
%------------------------------------------------------------------------------%
\def\disF{F}                  % distribution F
\def\disG{G}                  % distribution G
\def\disH{H}                  % distribution H
\def\disU{U}                  % distribution U
\def\disT{T}                  % distribution T
\def\SF{C^{\infty}(\G)}       % space of smooth functions
\def\TF{C_c^{\infty}(\G)}     % space of compactly supported smooth functions
\def\TFRd{C_c^{\infty}(\R^d)} % space of test functions
\def\SRd{\mathscr{S}(\R^d)}   % Schwartz space
\def\SR1{\mathscr{S}(\R)}     % Schwartz space
\def\B1#1{C^1(#1) }           % space of continuously differentiable functions
\def\Bm#1{C^m(#1) }           % space of continuously differentiable functions
\def\Ck{C^k(\G)}
\def\TD{\mathscr{S}'(\R^d)}   % space of temperated distributions
\def\E{\cal{E}}               % space of distributions with compact support
\def\D{\cal{D}}               % D
\def\FT{\mathcal{F}}          % Fourier transformation F
\def\FT{\mathscr{F}}          % Fourier transformation F

%------------------------------------------------------------------------------%
% other definitions                                                            %
%------------------------------------------------------------------------------%
\def\Proof{\textsc{Proof}}
\def\FRECHET{Fr\'{e}chet}
\def\ARZELA{Arzel\'{a}}
\def\GARDING{G\r{a}rding}
\def\HOELDER{H\"older}
\def\SCHROEDINGER{Schr\"odinger}
\def\JOERGENS{J\"orgens}
\def\WUEST{W\"uest}
\def\KAEHLER{K\"aehler}
\def\HOEFLINGER{H\"oflinger}
\def\POINCARE{Poincar\'{e}}
\def\FA{Functional Analysis}

\def\Author   {K.\,{\HOEFLINGER}}
\def\AuthorUC {K.\,HOEFLINGER}
\def\Version  {1.0}
\def\Email    {khoeflinger@t-online.de}
\def\khMath   {https://khoefli.github.io/khMath}
\def\Year     {2026}


\titleformat{\section}
{\normalfont \bfseries \filcenter}
{\thesection.\;}{0mm}{}

\titleformat{\subsection}%[runin]
{\normalfont \itshape}
{}{}{}

\titlespacing*{\section}   {0pt}{12pt}{4pt}
\titlespacing*{\subsection}{0pt}{6pt}{0pt}

%------------------------------------------------------------------------------%
%                                BEGIN OF DOCUMENT                             %
%------------------------------------------------------------------------------%
\begin{document}
\thispagestyle{empty}

\def\Title{Linear Operator Equations}

%\centerline{{\textit{Topics in Functional Analysis}}}
%\vspace{22pt}

\centerline{\large{\textbf{\Title}}}
\vspace{4pt}
\centerline{\small{{\textsc{\Author}}}}
\vspace{10pt}

\myfootnote { Latest Version: \today.}
\myfootnote { Email: \textit{{\Email}}, Url: \href{\khMath}{\khMath}.}


\fancyhead[C] {\small{\textsc{\Title}}}
\fancyhead[OL]{\small{\thepage}}
\fancyhead[ER]{\small{\thepage}}


%------------------------------------------------------------------------------%
%                                   CONTENT                                    %
%------------------------------------------------------------------------------%

\section{Introduction}
%------------------------------------------------------------------------------%
The study of so-called \textit{linear equations} is one of the main topics in
functional analysis, the mathematical theory generalizing linear algebra to
infinite-dimensional topological vector spaces.
This text serves as a gentle review of the fundamental functional analytic
methods to solve this type of equations, where we differ between
\begin{itemize}
\item[-]\textit{stationary equations} of the form $\linA x \eq y$,
\item[-]\textit{eigenvalue equations} of the form $\linA x - \lambda x \eq y$
        for arbitrary complex numbers $\lambda \xeof \C$, and finally
\item[-]\textit{evolutional equations} of the form
        $\dot{x}  + \linA x \eq y$ and
        $\ddot{x} + \linA x \eq y$,
        respectively.
\end{itemize}

The enumerated equations are supposed to be governed
by a \textit{linear operator} $\linA$
acting between complex Banach spaces $\banX$ and $\banY$,
called the \textit{domain space} and the \textit{range space} of $\linA$.

\subparagraph{}
It is often assumed that
the mapping $\linA$ is only defined
in a dense linear subspace $\domain{\linA} \xsubset \banX$,
called the \textit{domain of definition} of $\linA$,
so the operator under consideration appears as a mapping
\[
\linA \: \domain{\linA} \xsubset \banX \xto \banY
\]
fulfilling the linearity properties
\[
(L) \;
\begin{cases}[1.3]
\; \linA (x+y) = \linA x + \linA y,   \\
\; \linA(\lambda x) = \lambda \linA x \\
\end{cases}
\]
for all $x,y \xeof \domain{\linA}$ and $\lambda \xeof \C$.

\paragraph{}
In the special case $\banY \eq \banX$ we say that
$\linA$ is an operator \textit{acting in} $\banX$.
Furthermore,
good results in the analysis of the mentioned equations are obtained,
if $\banX \eq \hilH$ is a complex Hilbert space
with inner product $\scalar{x}{y} \xeof \C$
and norm $\norm{x} {:=} \, \scalar{x}{x}$ for $x,y \xeof \hilH$.

\section{The Basic Theory of Linear Operators}
%------------------------------------------------------------------------------%
Banach and Hilbert spaces can be seen as special classes of topological linear
spaces.
Hence they are endowed with a specific topology $\topT$,
called the \textit{norm topology} of such a space.
In fact, the collection $\topT$ of open sets is the coarsest topology
such that the norm mapping $x \mapsto \norm{x}$ remains continuous.

\paragraph{}
Given complex Banach spaces $\banX$ and $\banY$,
it makes sense to speak of \textit{continuous} linear operators between them.
From general topology it is known that a mapping $\phi\:\banX \xto \banY$
is continuous by definition,
if $\phi^{-1}(\openO)$ is open in $\banX$
for every open set $\openO \xsubset \banY$.

\subparagraph{}
Especially for linear operators the concept of continuity is well-defined.
It is strongly related to \textit{boundedness} of operators.
To be precise,
let $\linA$ be a linear operator, then $\linA$ is continuous \iff
$\linA$ is bounded in the sense that
\[
\norm{A} \eqdef \norm{\linA x}_{x \xeof \domain{\linA}, \norm{x} \eq 1} 
< \infty.
\]
In this case the value $\norm{A}$ is called the \textit{operator norm}
of $\linA$.
It is possible to extend the domain of definition of $\linA$ to the closure
of the linear subspace $\domain{\linA}$ in order to get a further bounded
linear operator $\linA_{ext}$ preserving the norm of $\linA$,
that is,
$\norm{A_{ext}} \eq \norm{A}$.

\paragraph{}
A linear operator $\compA \: \banX \xto \banY$ between Banach spaces
$\banX,\,\banY$ is called \textit{compact},
if the image of each bounded set under $\compA$ is relative compact,
or equivalently,
if $\{x_n\}_{n \in \N}$ is a bounded sequence in $\banX$,
then the sequence $\{\compA \, x_n\}_{n \in \N}$ in $\banY$
owns a convergent subsequence.

\subparagraph{}
For example,
a linear operator $\compA$ is compact,
if at least one of the spaces $\banX$ and $\banY$ is finite-dimensional.
In case of $\banY \eq \banX$,
the identity operator $\1_{\banX} \: \banX \xto \banX$ is compact
iff $\banX$ is finite-dimensional.
Compact operator own a couple of useful properties:

\begin{itemize}[leftmargin=9mm]
\item[(a)]
every compact linear operator is continuous;

\item[(b)]
the composition $\linA_1 \circ \linA_2$ of bounded linear operators
is compact,
if one of them is so;

\item[(c)]
the first and second property together imply that
the set $\CO{\banX}$ of compact linear operators $\compA\:\banX \xto \banX$
is a \textit{closed} ideal within the Banach algebra $\BO{\banX} $
of bounded linear operators $\linA \: \banX \xto \banX$,
hence it is complete and forms a Banach algebra in its own right;

\item[(d)]
finite linear combinations of compact linear operators are compact, too;

\item[(e)]
the limit of a norm convergent sequence of compact operators is compact.
\end{itemize}

\paragraph{} %- finite-rank operators
A linear operator $\linFR \: \banX \xto \banY$ is said
to have \textit{finite rank},
if the range $\range{\linFR}$ is finite-dimensional.
It is obvious that $\linFR$ is compact.

\subparagraph{}
Suppose that $(\linFR_n)$ is a sequence of finite-rank operators
converging to the operator $\linA$ {\wrt} the norm of $\BO{\banX,\banY}$.
Then $\linA$ is compact
and the set $\CO{\banX,\banY}$ of compact linear operators
proves to be a subset of the closure of the set $\FR{\banX,\banY}$
of finite-rank operators,
{\ie}
$\CO{\banX,\banY} \xsubset \overline{\FR{\banX,\banY}}$.

\subparagraph{}
Reversely,
in the setting of Hilbert spaces 
each compact linear operator $\compA$ appears in this way,
that is,
$\compA$ is the norm-limit of a sequence of finite-rank operators
(but this is not necessarily true in the general setting of Banach spaces).

\paragraph{}
We finally want to introduce a third class of Banach space operators,
defined by a special topological feature.

\subparagraph{}
A linear operator $\linA \: \domain{\linA} \xsubset \banX \xto \banY$
is called \textit{graph-closed} (or simply \textit{closed}) if the set
\[
\Gamma_{\linA} \eqdef
\{(x,\linA x) \: x \xeof \domain{\linA}\} \, \subset \, \banX \xtimes \banY,
\]
denoted the \textit{graph} of $\linA$,
is closed in the space $\banX \xtimes \banY$
{\wrt} the product topology $\topT_{\,\banX \xtimes \banY}$.
This property especially implies that
the kernel $\kernel{\linA}$ of $\linA$ forms a closed set in $\banX$.

\subparagraph{} %- sequentially closed operators
It is a matter of fact
that an operator $\linA$ is closed \iff it is \textit{sequentially closed},
that is,
$x_n \xto x$ and $\linA x_n \xto y$ implies
$x \xeof \domain{\linA}$ and $\linA x \eq y$,
for short
\[
\begin{rcases}
 x_n       \xto x \quad \\
 \linA x_n \xto y \quad \\
\end{rcases}
\quad \Longrightarrow \quad
x \xeof \domain{\linA} \, \and \, \linA x \eq y.
\]
This can easily proved by the fact
that a subset of a normed linear space is closed
\iff it is sequentially closed.
Checking the sequentially closeness is the usual method
to verify that a given linear operator is graph-closed.
But notice that the equivalence of these both notations is in general not valid,
if $\banX$ and $\banY$ are non-metrizable topological linear spaces.

\subparagraph{} %- graph norm
It is often helpful to make the domain of definition $\domain{\linA}$
of a linear operator $\linA$ into a normed linear space by its own.
This can be done by introducing the so-called \textit{graph norm}
\[
\norm{x}_{\linA} \eqdef \norm{x}_{\banX} \+ \norm{\linA x}_{\banY}
\]
for $x \xeof \domain{\linA}$.
In fact,
the pair $(\domain{\linA},\norm{\,\cdot\,}_{\linA})$
forms a normed linear space,
which constitutes a Banach space \iff 
$\linA$ is sequentially closed.
This equivalence provides a further method
to verify closeness of a given linear operator.

\paragraph{} %- closeness of operators based on Schauder estimates
The following theorem
uses the possible validity of an abstract \textit{Schauder} \textit{estimate}
to ensure sequentially closenss:

\begin{theorem}\label{BROWDER_THEOREM}
Let $\linA_0 \: \banX \xto \banY$ be a continuous linear operator,
where the Banach space $\banX$ is continuously and densely embedded
into the Banach space $\banY$ by a mapping $\iota \: \banX \xto \banY$.
If the estimate
\begin{equation}\label{SCHAUDER_ESTIMATE}
\norm{x}_{\banX} \, \leq \,
c \cdot
\{ \,\norm{\iota x}_{\banY} + \norm{\linA_0 \, x}_{\banY} \, \}
\end{equation}
is fulfilled for some $c {\,>\,} 0$ and all $x \xeof \banX$,
then the linear operator $\linA$ acting in $\banY$ with
$\domain{\linA} := \iota(\banX)$ and
$\linA\, \iota x := \linA_0 \, x$ for $x \xeof \banX$ is closed.
\end{theorem}

\paragraph{}
Let us emphasize that in general 
the concepts of closeness and continuity of a linear operator $\linA$
are not comparable.
However,
if $\banX$ and $\banY$ are Banach or Hilbert spaces
and $\domain{\linA}$ is a \textit{closed} set in $\banX$,
then closeness of $\linA$ is an easy consequence of continuity of $\linA$.
This is especially true in the following situation:
if $\linA$ is densely defined and continuous,
then $\linA$ can be extended to a continuous linear operator
$\linA_{ext} \: \banX \xto \banY$,
and since $\banX$ itself is closed,
the extension $\linA_{ext}$ is a closed operator.
Hence continuous linear operators of the form $\linA \: \banX \xto \banY$
fit in a natural way
into the class of densely defined and closed linear operators.

\subparagraph{} %-- the closed graph theorem
Reversely,
the \textit{closed graph theorem} claims that if $\domain{\linA}$ is closed,
then closeness implies continuity,
so both concepts are equivalent to each other
under this restrictive condition.

\paragraph{} %- bounded inverse theorem
The following application of the closed graph theorem is of outstanding meaning:
if a linear operator $\linA$ is closed and bijective,
then its inverse $\linA^{-1}$ is also linear and closed,
and since the domain of definition $\domain{\linA^{-1}}$
is the entire space $\banY$,
the closed graph theorem implies that $\linA^{-1} \: \banY \xto \banX$
is continuous and bounded,
respectively.
This fact is known as the \textit{bounded inverse theorem}
for closed linear operators.
But notice that the quotation is also valid for continuous linear operators
$\linA \: \banX \xto \banY$,
which is an immediate consequence
of \textsc{Banach}'s \textit{open mapping theorem}.

\paragraph{} %- summary on continuous and closed linear operators
In conclusion,
if the domain of definition $\domain{\linA}$ of a linear operator $\linA$
is a \textit{closed} subspace of the domain space $\banX$,
then the following chain of implications is valid:
\[
\linA \mbox{ compact }    \Longrightarrow
\linA \mbox{ bounded }    \Longleftrightarrow
\linA \mbox{ continuous } \Longleftrightarrow
\linA \mbox{ closed}. %  \Longrightarrow
%\linA \mbox{ closable}.
\]

\section{Stationary Equations}
%------------------------------------------------------------------------------%

\section{Spectral Theory and Eigenvalue Equations}
%------------------------------------------------------------------------------%

\section{Semigroups of Bounded Linear Operators}
%------------------------------------------------------------------------------%

\section{Evolution Equations of First Order}
%------------------------------------------------------------------------------%

\section{Evolution Equations of Second Order}
%------------------------------------------------------------------------------%

\section{The Spectral Decomposition Approach}
%------------------------------------------------------------------------------%

\section{Transformation Methods}
%------------------------------------------------------------------------------%

%------------------------------------------------------------------------------%
%                                 BIBLIOGRAPHY                                 %
%------------------------------------------------------------------------------%
\newpage
\thispagestyle{empty}

\vspace{8mm}
\begin{thebibliography}{99}
\begin{small}

\bibitem{EN}  K.-J.\,Engel, R.\,Nagel:
\textit{A Short Course on Operator Semigroups}.
Universitext, Springer, 2006.

\bibitem{Go}  J.\,Goldstein:
\textit{Semigroups of Linear Operators and Applications}.
Oxford University Press, 1985.

\bibitem{Ka} T.\,Kato:
\textit{Perturbation Theory for Linear Operators}, Springer Verlag.
Grundlehren der mathematischen Wissenschaften, Band 132, 1995.

\bibitem{Pa}  A.\,Pazy:
\textit{Semigroups of Linear Operators
        and Applications to Partial Differential Equations}.
Applied Mathematical Sciences Vol.\,44, Springer, 1983.

\bibitem{Yo}  K.\,Yosida:
\textit{Functional Analysis}.
Springer, 1980.

\end{small}
\end{thebibliography}
%------------------------------------------------------------------------------%
%                               END OF DOCUMENT                                %
%------------------------------------------------------------------------------%
\end{document}
